
\section{Folio 4}

\begin{ejercicio} % //TODO: hacer
    Sea $X$ un espacio de Banach y $T \in BL(X)$. Consideramos el polinomio:
    \[ P(t) = 1 + 2t + 3t^2 \]
    y el operador
    \[ P(T) = I + 2T + 3T^2 \]
    donde $I$ es la identidad. \\
    Demostrar las siguientes afirmaciones:
    \begin{itemize}
        \item[i)] si $\lambda$ es un valor propio de $T$, entonces $P(\lambda)$ es un valor propio de $P(T)$
        \item[ii)] $P(\sigma(T)) \subseteq \sigma(P(T))$ \qquad [$\sigma(\cdot)$ es el espectro].
    \end{itemize}
    \textit{Sugerencia: razonar por contradicción.}
\end{ejercicio}

\begin{ejercicio}
    Sea $X = C^1([0,1])$ y definimos en él la norma:
    \[ \|u\|_0 = |u(0)| + \|u'\|_\infty \qquad \forall u \in X, \]
    \begin{itemize}
        \item[i)] ¿el espacio normado $(X, \|u\|_0)$ es completo? (Motivar la respuesta)
        \item[ii)] Demostrar que $\|\cdot\|_0$ es una norma equivalente a la norma usual en $X$ definida por:
        \[ \|u\| = \|u\|_\infty + \|u'\|_\infty. \]
    \end{itemize}
    Para realizar la demostración de ambos enunciados, comenzaremos con el enunciado $(ii)$ para después usar este para demostrar el $(i)$. Antes de nada recordemos como viene dada la norma del infinito
    $$\|u\|_\infty = \max_{x \in [0,1]} |u(x)|\quad \forall u \in C^1([0,1])$$
    (se suele escribir $\sup$, pero como sabemos que es continua y en intervalo cerrado y acotado da igual)
    \begin{description}
        \item[$(ii)$] Recordemos que dos normas $\|\cdot\|_a$ y $\|\cdot\|_b$ en un espacio vectorial $X$ son equivalentes si existen dos constantes positivas $\alpha, \beta > 0$ tales que para cualquier elemento $u \in X$ se cumple la doble desigualdad
        $$\alpha \|u\|_b \le \|u\|_a \le \beta \|u\|_b$$
        Queremos ver que esto ocurre entre nuestra norma propuesta $\|u\|_0 = |u(0)| + \|u'\|_\infty$ y la norma usual $\|u\| = \|u\|_\infty + \|u'\|_\infty$.

        \begin{itemize}
            \item Para la \underline{primera desigualdad}, sabemos que para cualquier función continua $u \in C^1([0,1])$, el valor absoluto de la función evaluada en el punto $0$, $|u(0)|$, nunca puede superar al máximo valor absoluto que toma la función en todo el intervalo, el cual es por definición la norma infinito $\|u\|_\infty$. \\
            
            \noindent
            Por lo tanto, de forma directa
            $$|u(0)| \le \|u\|_\infty\implies |u(0)| + \|u'\|_\infty \le \|u\|_\infty + \|u'\|_\infty$$
            esto demuestra la primera parte de la equivalencia eligiendo como constante $\alpha = 1$
            $$\|u\|_0 \le 1 \cdot \|u\| $$

            \item Para la \underline{segunda desigualdad}, vamos a el TFC (Teorema Fundamental del Cálculo) para funciones derivables, que nos dice que para cualquier función $u \in C^1([0,1])$ y para cualquier punto $x \in [0,1]$, se cumple la siguiente igualdad
            $$\int_{0}^{x} u'(t) \, dt=u(x)-u(0)\implies u(x)=u(0)+\int_{0}^{x} u'(t) \, dt$$
            y veremos ahora las siguientes desigualdades tomando $x\in[0,1]$
            \begin{align*}
                |u(x)| \le |u(0)| + \left| \int_{0}^{x} u'(t) \, dt \right| &\le |u(0)| + \int_{0}^{x} |u'(t)| \, dt \overset{(*)}{\leq} |u(0)| + \int_{0}^{x} \|u'\|_\infty \, dt = \\ 
                &=|u(0)| + \|u'\|_\infty \cdot x \overset{(**)}{\leq} |u(0)| + \|u'\|_\infty=\|u\|_0
            \end{align*}
            donde en $(*)$ hemos usado que $|u'(t)| \le \|u'\|_\infty$ para todo $t \in [0,1]$, y en $(**)$ hemos usado que $x\in [0,1]$. Demostrada la desigualdad anterior tenemos lo siguiente
            \begin{align*}
                \|u\|_\infty = \sup_{x \in [0,1]} |u(x)| \le \|u\|_0 &\implies \|u\|_\infty + \|u'\|_\infty \le \|u\|_0 + \|u'\|_\infty\overset{(*)}{\leq} 2 \|u\|_0 \implies \\
                &\implies \dfrac{1}{2}\|u\| \le \|u\|_0
            \end{align*}
            donde en $(*)$ hemos usado que $\|u'\|_\infty \le |u(0)| + \|u'\|_\infty = \|u\|_0$, y tenemos finalmente la otra desigualdad buscada.
        \end{itemize}
        \item[$(i)$] Intuitivamente podemos ver que dadas dos normas $\|\cdot\|_a$ y $\|\cdot\|_b$ equivalentes en un espacio vectorial $X$. Entonces, el espacio normado $(X, \|\cdot\|_a)$ es completo si y solo si el espacio normado $(X, \|\cdot\|_b)$ es completo. \\
        
        \noindent
        Sabiendo ya que $X = C^1([0,1])$ dotado de su norma usual $\|u\| = \|u\|_\infty + \|u'\|_\infty$ es un espacio de Banach estándar (su completitud se hereda de que el espacio de funciones continuas $C([0,1])$ con la norma supremo es completo), entonces es claro que 
        \begin{equation*}
            (X, \|\cdot\|_0) \text{ es completo}
        \end{equation*}
    \end{description}
\end{ejercicio}


\begin{ejercicio} % //TODO: hacer
    Sea $\{F_n : \ell_2 \longrightarrow \mathbb{R}\}$ la sucesión de funcionales definidos por:
    \[ F_n(x) = x_{n+1} + x_{n+2} \qquad \text{si } x = (x_1, x_2, \dots, x_n, \dots) \in \ell_2. \]
    Demostrar que los funcionales $F_n$ son lineales y continuos y que $F_n$ converge $*$-débilmente a cero, pero $\{F_n\}$ no converge a cero en $(\ell_2)^*$.
\end{ejercicio}

\begin{ejercicio} 
    Sea $\{e_n\}_{n \in \mathbb{N}^+}$ un sistema ortonormal completo en un espacio de Hilbert. Sea $T$ el operador de $H$ en $H$ definido por:
    \[ T(x) = \alpha_1 e_2 + \alpha_2 e_3 + \dots + \alpha_n e_{n+1} + \dots \]
    para $x = \sum\limits_{n=1}^{+\infty} \alpha_n e_n$. Demostrar que $T$ es una isometría y encontrar su operador adjunto. \\

    \noindent
    Recordemos que un operador $T: H \longrightarrow H$ es una isometría si conserva la norma de todos los vectores del espacio, es decir, si se cumple
    \begin{equation*}
        \|T(x)\| = \|x\| \quad \forall x \in H
    \end{equation*}
    Dado que $\{e_n\}$ es un sistema ortonormal completo, cualquier vector $x \in H$ se puede escribir en términos de su serie de Fourier abstracta
    $$x = \sum_{n=1}^{+\infty} \alpha_n e_n=\sum_{n=1}^{+\infty} (x,e_n) e_n$$
    donde $\alpha_n = ( x, e_n )$ para cada $n$. Recordemos ahora la identidad de Parseval que dice

    \begin{equation*}
        \|x\|^2 = \sum_{n=1}^{+\infty} |(x, e_n)|^2 \quad \forall x\in \mathcal{H}
    \end{equation*}
    por lo que aplicandolo a nuestra imagen tendríamos
    \begin{equation*}
        \|T(x)\|^2 = \sum_{n=1}^{+\infty} |(T(x), e_n)|^2
    \end{equation*}
    Veámos el resultado para cada $n$:
    \begin{itemize}
        \item \tbf{Caso $n=1$}: tenemos
        \begin{equation*}
            (T(x), e_1) = \left( \alpha_1 e_2 + \alpha_2 e_3 + \dots, e_1 \right) = 0
        \end{equation*}
        porque $e_1$ es ortogonal a todos los demás $e_{n}$ con $n \geq 2$.
        \item \tbf{Caso $n \geq 2$}: tenemos
        \begin{equation*}
            (T(x), e_n) = \left( \alpha_1 e_2 + \alpha_2 e_3 + \dots +\alpha_{n-1} e_n + \dots, e_n \right) = \alpha_{n-1} = (x, e_{n-1})
        \end{equation*}
    \end{itemize} 
    por lo que tenemos 
    \begin{equation*}
        \|T(x)\|^2 = \sum_{n=1}^{+\infty} |(T(x), e_n)|^2 = |0|^2 + \sum_{n=2}^{+\infty} |(T(x), e_{n})|^2 = \sum_{n=2}^{+\infty} |\alpha_{n-1}|^2 = \sum_{n=1}^{+\infty} |\alpha_n|^2 = \|x\|^2
    \end{equation*}
    y se tiene demostrado lo que se quería. \\

    \noindent
    Por definición, el operador adjunto $T^*$ es el único operador lineal y acotado que satisface la identidad del producto escalar
    $$( T(x), y ) = ( x, T^*(y) ) \quad \forall x, y \in H$$
    Para deducir su fórmula, vamos a tomar dos vectores genéricos de nuestro espacio expresados en la base de Hilbert
    \begin{equation*}
        x = \sum_{n=1}^{+\infty} \alpha_n e_n, \quad y = \sum_{m=1}^{+\infty} \beta_m e_m
    \end{equation*}
    Desarrollemos el lado izquierdo de la identidad utilizando la definición de $T(x)$ y las propiedades de linealidad del producto escalar
    
    \begin{equation*}
        ( T(x), y ) = \left( \sum_{n=1}^{+\infty} \alpha_n e_{n+1}, \sum_{m=1}^{+\infty} \beta_m e_m \right) = \sum_{n=1}^{+\infty} \sum_{m=1}^{+\infty} \alpha_n \overline{\beta_m} ( e_{n+1}, e_m )\overset{(*)}{=} \sum_{n=1}^{+\infty} \alpha_n \overline{\beta_{n+1}}
    \end{equation*}
    donde en $(*)$ hemos usado que el producto interno es $0$ siempres excepto cuando $m=n+1$, y si desarrollamos un poco para igualarlo a un producto interno tendríamos
    \begin{equation*}
        ( T(x), y )=  \sum_{n=1}^{+\infty} \alpha_n \overline{\beta_{n+1}}=\left( \sum_{n=1}^{+\infty} \alpha_n e_n, \sum_{n=1}^{+\infty} \beta_{n+1} e_n \right)=\left(x, \sum_{n=1}^{+\infty} \beta_{n+1} e_n\right)
    \end{equation*}
    por lo que el operador adjunto $T^*$ viene definido de la siguiente manera

    \begin{equation*}
        T^*(y) = \beta_2 e_1 + \beta_3 e_2 + \dots + \beta_{n+1} e_n + \dots \quad \text{para } y = \sum_{n=1}^{+\infty} \beta_n e_n
    \end{equation*}
\end{ejercicio}

\begin{ejercicio}
    Sea el siguiente operador
    \Func{T}{\ell_2}{\ell_2}{(x_1, \dots, x_n, \dots)}{(x_1, x_2, 0, 0, x_4, x_5, \dots, x_n, \dots) }
    Determinar el espectro de $T$. Identificar los posibles valores propios y los correspondientes espacios propios. ¿Es el operador $T$ compacto? \\
    \noindent
    Antes de nada es importante notar que $T$ es trivialmente lineal y para la continuidad se ve tan fácil como sigue sabiendo la equivalencia con la acotación
    $$\|T(x)\|_2^2 = |x_1|^2 + |x_2|^2 + 0^2 + 0^2 + |x_4|^2 + |x_5|^2 + \dots \leq \sum_{n=1}^{\infty} |x_n|^2 = \|x\|_2^2$$
    tomando la raíz, tenemos $\|T(x)\|_2 \leq 1 \cdot \|x\|_2 \implies \|T\| \leq 1$. El operador es acotado (continuo). \\ \\
    \noindent
    Pensemos en lo que hace $T$. Deja fijas infinitas componentes ($x_1, x_2, x_4, x_5, \dots$). Los operadores compactos en dimensión infinita tienen que ``aplastar'' o encoger las infinitas dimensiones hacia el cero (como el $\frac{1}{n}$ o el $\frac{1}{2^n}$ de los ejercicios anteriores). Si un operador deja intactas infinitas dimensiones, no puede ser compacto, por lo que nuestra idea será demostrar que NO es compacto. \\

    \noindent
    Supongamos por reducción al absurdo que $T$ es compacto. Consideremos la sucesión de vectores de la base canónica $\{e_n\}_{n \geq 5} \subseteq \ell_2$, donde $e_n = (0, \dots, 0, 1, 0, \dots)$ con el $1$ en la posición $n$. Sabemos que:
    \begin{itemize}
        \item Esta sucesión está acotada porque $\|e_n\|_2 = 1$ para todo $n$.
        \item Si aplicamos el operador a partir de $n \geq 5$, el operador los deja completamente idénticos: $T(e_n) = e_n$.
    \end{itemize}
    Por definición de operador compacto, si una sucesión está acotada, su imagen $\{T(e_n)\}$ debe contener una subsucesión convergente. En nuestro caso, la sucesión de imágenes es la propia $\{e_n\}_{n \geq 5}$. Veamos la distancia entre dos elementos cualesquiera de esta sucesión de imágenes ($n \neq m$ con $n, m \geq 5$)
    $$\|T(e_n) - T(e_m)\|_2^2 = \|e_n - e_m\|_2^2 = 1^2 + (-1)^2 = 2 \implies \|T(e_n) - T(e_m)\|_2 = \sqrt{2}$$
    como la distancia entre cualquier par de términos es siempre una constante fija ($\sqrt{2}$), es imposible extraer ninguna subsucesión que sea de Cauchy y que pueda converger, por lo que $T$ \underline{no es compacto}. \\

    \noindent
    Ahora estudiaremos los \underline{valores propios} y los \underline{subespacios propios asociados}. Para encontrar los valores propios, planteamos la ecuación de autovalores $T(x) = \lambda x$ para un vector no nulo $x \neq 0$. Lo escribimos componente a componente para ver qué pasa:
    \begin{equation*}
        \begin{cases}
            x_1 = \lambda x_1 \implies (1 - \lambda)x_1 = 0 \\
            x_2 = \lambda x_2 \implies (1 - \lambda)x_2 = 0 \\
            0 = \lambda x_3 \implies \lambda x_3 = 0 \\
            0 = \lambda x_4 \implies \lambda x_4 = 0 \\
            x_{n-1} = \lambda x_n \quad \text{para } n \geq 5
        \end{cases}
    \end{equation*}
    Analizando este sistema tan sencillo, vemos que solo hay dos números candidatos que pueden generar soluciones no nulas:
    \begin{itemize}
        \item \tbf{Caso $\lambda=1$}: sustituyendo $\lambda=1$ en las ecuaciones, obtenemos que $0 \cdot x_1 = 0$ y $0 \cdot x_2 = 0$, lo que significa que $x_1$ y $x_2$ pueden ser cualquier número. Sin embargo, las ecuaciones para $x_3$ y $x_4$ se convierten en $0 = 0$, lo que significa que $x_3$ y $x_4$ deben ser nulos.  \\
        \noindent
        Por último, de la ecuación (5) para $n=5$, tenemos que $x_4 = 1 \cdot x_5$. Como sabemos que $x_4 = 0$, entonces obligatoriamente $x_5 = 0$, y así ocurre sucesivamente para $n\geq 6$, obteniendo por tanto, que todos deben ser nulos excepto $x_1$ y $x_2$, es decir, el subespacio propio asociado es el siguiente
        $$V_1 = \ker(T - I) = \text{span}\{e_1, e_2\}$$

        \item \tbf{Caso $\lambda=0$}: sustituyendo $\lambda=0$ en las ecuaciones, obtenemos que $x_1 = 0$, $x_2 = 0$, $x_n = 0$ para todo $n \geq 4$. Sin embargo, las ecuaciones para $x_3$ y $x_4$ se convierten en $0 = 0$, lo que significa que $x_3$ y $x_4$ pueden ser cualquier número. De la ecuación (5) sustituyendo $\lambda = 0$, tenemos que $x_{n-1} = 0 \cdot x_n \implies x_{n-1} = 0 \quad \text{para } n \geq 5$. Esto significa que $x_4 = 0, x_5 = 0, x_6 = 0, \dots$, aqui se obliga a $x_4=0$, pero $x_3$ no se obliga a ser cero, por lo que el espacio propio asociado a $\lambda = 0$ es el siguiente
        $$V_0 = \ker(T) = \text{span}\{e_3\}$$

        \item \tbf{Caso $\lambda \neq 1$ y $\lambda \neq 0$}: sustituyendo estos valores en las ecuaciones, obtenemos que $x_1 = 0$, $x_2 = 0$, $x_3 = 0$, $x_4 = 0$, y $x_n = 0$ para todo $n \geq 5$. Por tanto, la única solución es el vector nulo, lo que significa que no hay más valores propios.
    \end{itemize}
    Los valores propios son exactamente $e(T) = \{0, 1\}$. \\

    \noindent
    Estudiaremos ahora el \underline{espectro} de $T$. Tomemos ahora $\lambda \neq 0,1$ y vamos a intentar despejar de forma única el vector $x$ en función de $y$, recordando primero que la ecuación general es $(T - \lambda I)x = y$, es decir, que
    \begin{equation*}
        \left( \left(1 - \lambda\right)x_1, \left(1 - \lambda\right)x_2, -\lambda x_3, -\lambda x_4, x_4-\lambda x_5 , \dots, x_{n-1}-\lambda x_n , \dots \right) = (y_1, y_2, y_3, y_4, y_5, y_6, \dots)
    \end{equation*}
    Despejando cada componente $x_n$ de forma única, obtenemos:
    \begin{itemize}
        \item $x_1 = \dfrac{y_1}{1 - \lambda}$
        \item $x_2 = \dfrac{y_2}{1 - \lambda}$
        \item $x_3 = \dfrac{y_3}{-\lambda}$
        \item $x_4 = \dfrac{y_4}{-\lambda}$
        \item $x_n = \dfrac{x_{n-1} - y_n}{\lambda} \implies x_n = -\dfrac{y_4}{\lambda^{n-3}} - \sum\limits_{k=5}^{n} \dfrac{y_k}{\lambda^{n-k+1}} \quad \text{para } n \geq 5$
    \end{itemize}
    como $\lambda \neq 0,1$, este despeje sigue siendo completamente único componente a componente (garantiza la inyectividad). Ahora tenemos que ver si exactamente ese $x$ despejado pertenece a $\ell_2$ para cada $y \in \ell_2$, para ello veremos dos casos
    \begin{itemize}
        \item \tbf{Caso $|\lambda |\leq 1$}: para demostrar que el operador no es sobreyectivo, basta con encontrar un único vector $y \in \ell_2$ cuya solución $x$ quede fuera de $\ell_2$. Tomamos el vector de la base canónica 
        $$y = e_4 = (0, 0, 0, 1, 0, 0, \dots) \in \ell_2$$
        Sustituyendo sus componentes ($y_4 = 1$ y el resto $0$) en nuestra fórmula general de la solución para $n \ge 4$
        $$x_n = -\frac{1}{\lambda^{n-3}}$$
        Ahora veremos si $x$ pertenece a $\ell_2$ o no, para ello evaluamos su norma al cuadrado
        $$\|x\|_2^2 = \sum_{n=4}^{\infty} |x_n|^2 = \sum_{n=4}^{\infty} \frac{1}{|\lambda|^{2(n-3)}}= \sum_{m=1}^{\infty} \left(\frac{1}{|\lambda|^2}\right)^{m}$$
        Como $|\lambda| \leq 1$, entonces $\frac{1}{|\lambda|^2} \geq 1$, por lo que la serie no converge, lo que significa que $x \notin \ell_2$. Por tanto, el operador no es sobreyectivo, lo que implica que $\lambda$ pertenece al espectro $\sigma(T)$.
        \item \tbf{Caso $|\lambda| > 1$}: es más complejo.
    \end{itemize}
    Lo que haremos es sabiendo que $\ell_2$ es espacio de Hilbert, aplicaremos que sabemos que 
    \begin{equation*}
        \sigma(T)  \subseteq \{\lambda \in \mathbb{C} : |\lambda| \leq \|T\|\}
    \end{equation*}
    y con el caso $|\lambda| \leq 1$ ya quedaría demostrada la igualdad del espectro. Para ello veámos la norma del operador (recordemos que al principio ya demostramos que $\|T\|\leq 1$). Para probar que la norma es exactamente $1$, basta con testear un vector con norma 1 que no se altere, por ejemplo el primer vector canónico $e_1 = (1, 0, 0, \dots)$. Como $\|e_1\|_2 = 1$ y $T(e_1) = e_1$, tenemos que $\|T(e_1)\|_2 = 1$. Por definición de norma de un operador, concluimos que:
    $$\|T\| = \sup_{x \neq 0} \frac{\|T(x)\|_2}{\|x\|_2} \geq \frac{\|T(e_1)\|_2}{\|e_1\|_2} = 1$$
    y se tiene finalmente que
    \begin{equation*}
        \sigma(T) = \{\lambda \in \mathbb{C} : |\lambda| \leq 1\}
    \end{equation*}
\end{ejercicio}

\begin{ejercicio}
    Sea $X$ un espacio de Banach y $T \in BL(X)$. Demostrar que $\lambda \in \rho(T)$ (resolvente de $T$) si y solo si:
    \begin{itemize}
        \item[(i)] $(T - \lambda I)(X)$ es denso en $X$
        \item[(ii)] $\exists \, \alpha > 0 : \quad \|(T-\lambda I)(x)\|=\|T(x) - \lambda x\| \geq \alpha \|x\| \qquad \forall x \in X$
    \end{itemize}

    \noindent
    Para simplificar la escritura en toda la demostración y que la notación no nos emborrache, vamos a definir el operador intermedio
    $$A = T - \lambda I$$
    donde claramente como $T \in BL(X)$ e $I \in BL(X)$, su resta también es un operador lineal y acotado ($A \in BL(X)$). Con esta notación lo que se nos pide demostrar es lo siguiente

    $$\lambda \in \rho(T) \iff A(X) \text{ es denso en } X \quad \text{y} \quad \exists \alpha> 0 : \quad \|Ax\| \geq \alpha\|x\| \quad \forall x \in X$$
    Lo veremos por doble implicación.
    \begin{description}
        \item[$\Longrightarrow)$] Sea $\lambda \in \rho(T)$, entonces $A$ es biyectiva, veámos como probar $(i)$ y $(ii)$:
        \begin{itemize}
            \item[$(i)$] Como $A$ es sobreyectiva, entonces $A(X) =\overline{A(X)}= X$, lo que implica que $A(X)$ es denso en $X$.
            \item[$(ii)$] Como $A$ biyectiva, existe su inversa
            \begin{equation*}
                A^{-1} : X \longrightarrow X
            \end{equation*}
            y por teorema de la aplicación abierta sabemos que $A^{-1}$ es continua, lo que implica que está acotada, es decir, tenemos
            $$\|A^{-1}y\| \leq \|A^{-1}\| \cdot \|y\|\quad \forall y \in X$$
            donde $\|A^{-1}\|>0$. Haciendo el cambio $y = Ax$ (lo cual es válido porque $A$ es biyectiva), tenemos $A^{-1}y = x$. Sustituyendo
            $$\|x\| \leq \|A^{-1}\| \cdot \|Ax\| \implies \|Ax\| \geq \frac{1}{\|A^{-1}\|} \|x\|$$
            Definiendo $\alpha = \frac{1}{\|A^{-1}\|} > 0$, obtenemos la condición (ii): $\|Ax\| \geq \alpha \|x\|$.
        \end{itemize}
        \item[$\Longleftarrow)$] Supongamos que se cumplen las condiciones (i) y (ii). Para demostrar que $\lambda \in \rho(T)$, es decir, que $A$ es biyectiva, tenemos que demostrar que $A$ es inyectiva y sobreyectiva.
        \begin{itemize}
            \item \tbf{Inyectividad}: queremos ver que si $Ax = 0 \implies x = 0$ (por ser $A$ lineal). Usamos la condición (ii)
            $$0 = \|0\| = \|Ax\| \geq \alpha \|x\|$$
            Como $\alpha > 0$, la única forma de que $\alpha \|x\| \leq 0$ es que $\|x\| = 0$, lo que significa que $x = 0$. Por tanto, $A$ es inyectiva.
            \item \tbf{Sobreyectividad}: lo que primero probaremos aqui es que la imagen es un conjunto cerrado, es decir, $A(X)=\overline{A(X)}$, puesto que entonces tomando $(i)$, ya sabemos que es sobreyectiva, puesto que por ser denso en $X$, tenemos
            $$ A(X)=\overline{A(X)}=X$$
            Tomamos una sucesión de elementos en la imagen, $y_n = Ax_n \in A(X)$, tal que $y_n \to y$ (en norma) en $X$. Queremos ver que el límite $y$ pertenece a $A(X)$. Como $y_n$ converge, es una sucesión de Cauchy, por lo que 
            $$\|y_n - y_m\| \to 0 \quad \text{ cuando } n,m \to \infty$$
            Usando la condición (ii) y la linealidad de $A$
            $$\|y_n - y_m\| = \|Ax_n - Ax_m\| = \|A(x_n - x_m)\| \geq \alpha \|x_n - x_m\|$$
            entonces tenemos que
            \begin{equation*}
                \|y_n - y_m\| \to 0 \implies \alpha \|x_n - x_m\| \to 0  \overset{(*)}{\implies} \|x_n - x_m\| \to 0
            \end{equation*}
            donde en $(*)$ hemos usado que $\alpha > 0$. Esto significa que $\{x_n\}$ es una sucesión de Cauchy en $X$, y como $X$ completo entonces $\{x_n\}$ converge para algún elemento de $X$, es decir
            $$x_n \to x_0 \quad \text{ donde }x_0 \in X$$ 
            Al ser $A$ un operador continuo, preserva los límites en norma 
            $$y_n = Ax_n \to Ax_0$$ 
            y por unicidad del límite, $y = Ax_0$. Como $x_0 \in X$, concluimos que $y \in A(X)$, por lo que $A(X)$ es cerrado, y por lo explicado al inicio, $A$ sobreyectiva.
        \end{itemize}
    \end{description}
\end{ejercicio}

\begin{ejercicio}
    Sea $X$ un espacio normado. Definimos a $\{x_n\}$ como sucesión débilmente de Cauchy en $X$ si 
    \begin{equation*}
        \{f(x_n)\} \text{ es una sucesión de Cauchy en } \mathbb{R} \quad \forall f \in X^* 
    \end{equation*}
    Demostrar que toda sucesión débilmente de Cauchy está acotada (en norma) en $X$. \\
    
    \noindent
    Sea $f \in X^*$ un funcional cualquiera fijo. Por hipótesis, la sucesión numérica $\{f(x_n)\} \subset \mathbb{R}$ es una sucesión de Cauchy. Dado que toda sucesión de Cauchy en $\mathbb{R}$ está acotada, deducimos que para cada funcional $f$ existe una constante real $C_f > 0$ (que depende de $f$) tal que
    \begin{equation}\label{eq:acotacion}
        |f(x_n)| \le C_f \quad \forall n \in \mathbb{N} 
    \end{equation}
    Para ver más claramente que efectivamente toda sucesión de Cauchy está acotada en $\R$, sabemos primero la definición de dicha sucesión como sigue
    \begin{equation*}
        \forall \varepsilon > 0, \ \exists N \in \mathbb{N} \ : \ |a_n - a_m| < \varepsilon \quad \forall n, m \ge N
    \end{equation*}
    y si tomamos $\varepsilon=1$, sabemos que existe un subíndice concreto $N \in \mathbb{N}$ tal que
    $$|a_n - a_m| < 1 \quad \forall n, m \ge N$$
    donde tomando $m=N$, tenemos
    \begin{equation*}
        |a_n - a_N| < 1 \quad \forall n \ge N \implies |a_n| \le |a_n - a_N| + |a_N| < 1 + |a_N| \quad \forall n \ge N
    \end{equation*}
    entonces tomamos
    \begin{equation*}
        M = \max\big\{ |a_1|, |a_2|, \dots, |a_{N-1}|, 1 + |a_N| \big\} \implies |a_n| \le M \quad \forall n \in \mathbb{N}
    \end{equation*}
    Visto este pequeño detalle seguimos con la resolución del ejericicio. Para cada elemento $x_n$ de nuestra sucesión, consideramos su imagen mediante la inyección canónica $J: X \to X^{**}$. Definimos los funcionales $\xi_n \in X^{**}$ como
    $$ \xi_n(f) = J(x_n)(f) = f(x_n) \quad \forall f \in X^* $$
    Recordemos que la inyección canónica es una isometría, lo que significa que la norma de $\xi_n$ en el bidual es exactamente igual a la norma de $x_n$ en el espacio original
    $$ \|\xi_n\|_{X^{**}} = \|x_n\|_X$$
    Como $X$ espacio normado y $\R$ espacio de Banach, entonces $X^{**}$ es espacio de Banach. Mirando la familida de operadores lineales y continuos $\{\xi_n\}_{n \in \mathbb{N}}$, evaluemos esta familia de operadores de forma puntual para un elemento $f \in X^*$ cualquiera
    $$ \sup_{n \in \mathbb{N}} |\xi_n(f)| = \sup_{n \in \mathbb{N}} |f(x_n)| $$
    y como dijimos en la Ecaución (\ref{eq:acotacion}), tenemos
    $$ \sup_{n \in \mathbb{N}} |\xi_n(f)| \le C_f < +\infty \quad \forall f \in X^* $$
    lo que significa que la familia de operadores $\{\xi_n\}$ está acotada puntualmente sobre el espacio de Banach $X^*$. Al cumplirse las hipótesis del Principio de Acotación Uniforme (PAU), este teorema nos garantiza que la familia también debe estar acotada uniformemente en norma, es decir
    \begin{equation*}
        \exists M>0 : \quad \|\xi_n\|_{X^{**}} \le M \quad \forall n \in \mathbb{N}\implies \|x_n\|_X= \|\xi_n\|_{X^{**}}  \le M \quad \forall n \in \mathbb{N}
    \end{equation*}
    lo que demuestra formalmente que la sucesión $\{x_n\}$ está acotada en norma en el espacio $X$.
\end{ejercicio}

